\section{Introduction}\label{sec:intro}

The category $\oCat$ of strict, globular $\omega$-categories bears a
``canonical'' model structure structure introduced
in~\cite{lafontetal:folkms}. In this structure, the cofibrant objects
are of the form $\frcat P$ where $P$ is a polygraph
(see~\cite{abgmmm:polybk}) and $\frcat P$ denotes the \oo-category
freely generated by $P$. Let $\Chp$ be the category of chain complexes
of abelian groups in non-negative degree equipped with its usual
projective model structure. There is a left-Quillen abelianization
functor $\abel:\oCat\to \Chp$ (see~\ref{subsec:abel} below). Let then $\LL\abel:\Loc{\oCat}\to\Loc{\Chp}$ denote the
left-derived functor of $\abel$ between the corresponding homotopy
categories. The {\em polygraphic homology} of an \oo-category $C$
is by definition
\[\Hopol(C)=\LL\abel(C).\]
To compute this homology, we first produce
a cofibrant replacement of $C$, namely a polygraph $P$ together with an 
an acyclic fibration $p:\frcat P\to C$, then apply the functor $\abel$
directly to $\frcat P$, whence the name.

Now for each integer $k\geq 0$, an \ook k-category is an \oo-category
in which all cells are invertible in dimensions $>k$. The full
subcategory of $\oCat$ whose objects are \ook k-categories is denoted
by $\opCat k$. For instance $\opCat 0$ is just the category of strict
\oo-groupoids.
The canonical inclusion functor $\oemb k:\opCat k\to\oCat$ admits a left
adjoint $\opifun k:\oCat\to\opCat k$. It turns out that the canonical model 
structure on $\oCat$ transfers to $\opCat k$ via $\oemb k$ and that
$\opifun k$ is again left-Quillen. On the other hand, the functor
$\lambda$ easily factors through $\opifun k$. Thus we get an
abelianization functor $\abelk k:\opCat k\to\Chp$ such that
$\abel=\abelk k\circ\opifun k$, and $\LL\abel$ is naturally isomorphic
to $\LL\abelk k\circ \LL\opifun k$. This leads to define the
polygraphic homology of an \ook k-category $C$ by
\[\Hop{k}(C)=\LL\abelk k(C).\]

 
  